%%%%%%%%%%%%%%%%%%%%%%%%%%%%%%%%%%%%%%%%%%%%%%%%%%%%%%%%%%%%
% 5.8 COMPLEX ANALYSIS
% The Composition of Advanced Mathematics
%%%%%%%%%%%%%%%%%%%%%%%%%%%%%%%%%%%%%%%%%%%%%%%%%%%%%%%%%%%%

\section{Complex Analysis}
\label{sec:complex-analysis}

\prerequisite{
Complex numbers, functions, limits, continuity, differentiation,
integration, sequences and series, power series, multivariable calculus,
and basic real analysis.
}

\learningobjectives{
By the end of this section, the reader should be able to:
\begin{itemize}
    \item represent complex numbers algebraically, geometrically, and in polar form;
    \item use modulus, argument, conjugation, and Euler's formula;
    \item analyze limits and continuity of complex-valued functions;
    \item define complex differentiability and holomorphic functions;
    \item apply the Cauchy--Riemann equations;
    \item recognize harmonic functions and harmonic conjugates;
    \item work with the complex exponential, logarithm, powers, and roots;
    \item understand branches and branch cuts;
    \item evaluate complex contour integrals;
    \item apply Cauchy's Integral Theorem and Cauchy's Integral Formula;
    \item understand analyticity and power-series representations;
    \item work with zeros and isolated singularities;
    \item construct Laurent series;
    \item compute residues;
    \item apply the Residue Theorem;
    \item evaluate selected real integrals using residues;
    \item apply the Maximum Modulus Principle, Liouville's Theorem,
          and the Fundamental Theorem of Algebra;
    \item understand the Argument Principle and Rouch\'e's Theorem;
    \item recognize conformal mappings and M\"obius transformations.
\end{itemize}
}

%%%%%%%%%%%%%%%%%%%%%%%%%%%%%%%%%%%%%%%%%%%%%%%%%%%%%%%%%%%%
\subsection{The Complex Plane}
\label{subsec:complex-plane}
%%%%%%%%%%%%%%%%%%%%%%%%%%%%%%%%%%%%%%%%%%%%%%%%%%%%%%%%%%%%

A complex number has the form

\[
\boxed{
z=x+iy,
}
\]

where

\[
x,y\in\R
\]

and

\[
\boxed{
i^2=-1.
}
\]

The number \(i\) is called the \emph{imaginary unit}.

The set of complex numbers is denoted

\[
\boxed{
\mathbb{C}.
}
\]

A complex number may be identified with the point

\[
(x,y)\in\R^2.
\]

Thus the complex plane provides a geometric representation of
\(\mathbb{C}\).

%%%%%%%%%%%%%%%%%%%%%%%%%%%%%%%%%%%%%%%%%%%%%%%%%%%%%%%%%%%%
\subsection{Real and Imaginary Parts}
\label{subsec:real-imaginary-parts}
%%%%%%%%%%%%%%%%%%%%%%%%%%%%%%%%%%%%%%%%%%%%%%%%%%%%%%%%%%%%

If

\[
z=x+iy,
\]

then

\[
\boxed{
\operatorname{Re}(z)=x
}
\]

and

\[
\boxed{
\operatorname{Im}(z)=y.
}
\]

These are called the real and imaginary parts of \(z\).

Two complex numbers are equal precisely when their real and imaginary
parts agree.

%%%%%%%%%%%%%%%%%%%%%%%%%%%%%%%%%%%%%%%%%%%%%%%%%%%%%%%%%%%%
\subsection{Complex Conjugation}
\label{subsec:complex-conjugation}
%%%%%%%%%%%%%%%%%%%%%%%%%%%%%%%%%%%%%%%%%%%%%%%%%%%%%%%%%%%%

\begin{definitionbox}{Complex Conjugate}
If

\[
z=x+iy,
\]

then the complex conjugate of \(z\) is

\[
\boxed{
\overline{z}=x-iy.
}
\]
\end{definitionbox}

Important identities include

\[
\boxed{
z+\overline{z}=2\operatorname{Re}(z),
}
\]

\[
\boxed{
z-\overline{z}=2i\operatorname{Im}(z),
}
\]

and

\[
\boxed{
z\overline{z}=x^2+y^2.
}
\]

Also,

\[
\overline{zw}
=
\overline{z}\,\overline{w},
\]

and

\[
\overline{z+w}
=
\overline{z}+\overline{w}.
\]

%%%%%%%%%%%%%%%%%%%%%%%%%%%%%%%%%%%%%%%%%%%%%%%%%%%%%%%%%%%%
\subsection{Modulus}
\label{subsec:complex-modulus}
%%%%%%%%%%%%%%%%%%%%%%%%%%%%%%%%%%%%%%%%%%%%%%%%%%%%%%%%%%%%

\begin{definitionbox}{Complex Modulus}
For

\[
z=x+iy,
\]

the modulus of \(z\) is

\[
\boxed{
|z|
=
\sqrt{x^2+y^2}.
}
\]
\end{definitionbox}

Geometrically, \(|z|\) is the distance from \(z\) to the origin in the
complex plane.

Since

\[
z\overline{z}=|z|^2,
\]

we have

\[
\boxed{
|z|
=
\sqrt{z\overline{z}}.
}
\]

%%%%%%%%%%%%%%%%%%%%%%%%%%%%%%%%%%%%%%%%%%%%%%%%%%%%%%%%%%%%
\subsection{Properties of the Modulus}
\label{subsec:properties-complex-modulus}
%%%%%%%%%%%%%%%%%%%%%%%%%%%%%%%%%%%%%%%%%%%%%%%%%%%%%%%%%%%%

For \(z,w\in\mathbb{C}\),

\[
\boxed{
|zw|=|z||w|,
}
\]

\[
\boxed{
\left|\frac{z}{w}\right|
=
\frac{|z|}{|w|}
\qquad (w\neq0),
}
\]

and

\[
\boxed{
|z+w|
\leq
|z|+|w|.
}
\]

The final inequality is the triangle inequality.

The reverse triangle inequality is

\[
\boxed{
\bigl||z|-|w|\bigr|
\leq
|z-w|.
}
\]

%%%%%%%%%%%%%%%%%%%%%%%%%%%%%%%%%%%%%%%%%%%%%%%%%%%%%%%%%%%%
\subsection{Argument}
\label{subsec:complex-argument}
%%%%%%%%%%%%%%%%%%%%%%%%%%%%%%%%%%%%%%%%%%%%%%%%%%%%%%%%%%%%

For \(z\neq0\), an argument of \(z\) is an angle \(\theta\) satisfying

\[
z
=
|z|
(\cos\theta+i\sin\theta).
\]

Because angles differing by \(2\pi\) represent the same direction,

\[
\boxed{
\arg z
=
\theta+2\pi k,
\qquad
k\in\Z.
}
\]

Thus the argument is naturally multivalued.

A selected single-valued argument is called a principal argument and is
often denoted

\[
\boxed{
\operatorname{Arg}z.
}
\]

One common convention is

\[
-\pi<\operatorname{Arg}z\leq\pi.
\]

%%%%%%%%%%%%%%%%%%%%%%%%%%%%%%%%%%%%%%%%%%%%%%%%%%%%%%%%%%%%
\subsection{Polar Form}
\label{subsec:complex-polar-form}
%%%%%%%%%%%%%%%%%%%%%%%%%%%%%%%%%%%%%%%%%%%%%%%%%%%%%%%%%%%%

Every nonzero complex number can be written

\[
\boxed{
z
=
r(\cos\theta+i\sin\theta),
}
\]

where

\[
r=|z|>0.
\]

This is the polar form of \(z\).

Multiplication becomes particularly simple:

\[
z_1z_2
=
r_1r_2
\left[
\cos(\theta_1+\theta_2)
+
i\sin(\theta_1+\theta_2)
\right].
\]

Thus multiplication multiplies moduli and adds arguments.

%%%%%%%%%%%%%%%%%%%%%%%%%%%%%%%%%%%%%%%%%%%%%%%%%%%%%%%%%%%%
\subsection{Euler's Formula}
\label{subsec:eulers-formula-complex}
%%%%%%%%%%%%%%%%%%%%%%%%%%%%%%%%%%%%%%%%%%%%%%%%%%%%%%%%%%%%

One of the fundamental identities of complex analysis is

\[
\boxed{
e^{i\theta}
=
\cos\theta+i\sin\theta.
}
\]

Therefore polar form becomes

\[
\boxed{
z=re^{i\theta}.
}
\]

This gives

\[
\boxed{
e^{i\pi}+1=0,
}
\]

Euler's famous identity.

%%%%%%%%%%%%%%%%%%%%%%%%%%%%%%%%%%%%%%%%%%%%%%%%%%%%%%%%%%%%
\subsection{Powers and De Moivre's Formula}
\label{subsec:de-moivre}
%%%%%%%%%%%%%%%%%%%%%%%%%%%%%%%%%%%%%%%%%%%%%%%%%%%%%%%%%%%%

If

\[
z=re^{i\theta},
\]

then for an integer \(n\),

\[
\boxed{
z^n
=
r^ne^{in\theta}.
}
\]

Equivalently,

\[
\boxed{
(\cos\theta+i\sin\theta)^n
=
\cos(n\theta)+i\sin(n\theta).
}
\]

This is De Moivre's formula.

%%%%%%%%%%%%%%%%%%%%%%%%%%%%%%%%%%%%%%%%%%%%%%%%%%%%%%%%%%%%
\subsection{Complex Roots}
\label{subsec:complex-roots}
%%%%%%%%%%%%%%%%%%%%%%%%%%%%%%%%%%%%%%%%%%%%%%%%%%%%%%%%%%%%

To solve

\[
z^n=w,
\]

write

\[
w=re^{i\theta}.
\]

The \(n\) distinct roots are

\[
\boxed{
z_k
=
r^{1/n}
e^{i(\theta+2\pi k)/n},
\qquad
k=0,1,\ldots,n-1.
}
\]

Thus the roots are equally spaced around a circle.

For

\[
z^n=1,
\]

the roots are

\[
\boxed{
z_k
=
e^{2\pi ik/n}.
}
\]

These are the \(n\)th roots of unity.

%%%%%%%%%%%%%%%%%%%%%%%%%%%%%%%%%%%%%%%%%%%%%%%%%%%%%%%%%%%%
\subsection{Complex Functions}
\label{subsec:complex-functions}
%%%%%%%%%%%%%%%%%%%%%%%%%%%%%%%%%%%%%%%%%%%%%%%%%%%%%%%%%%%%

A complex-valued function may be written

\[
\boxed{
f(z)=u(x,y)+iv(x,y),
}
\]

where

\[
z=x+iy
\]

and

\[
u,v:\R^2\to\R.
\]

The functions \(u\) and \(v\) are respectively the real and imaginary
parts of \(f\).

Complex analysis therefore combines ideas from real multivariable
calculus with the algebraic structure of \(\mathbb{C}\).

%%%%%%%%%%%%%%%%%%%%%%%%%%%%%%%%%%%%%%%%%%%%%%%%%%%%%%%%%%%%
\subsection{Complex Limits}
\label{subsec:complex-limits}
%%%%%%%%%%%%%%%%%%%%%%%%%%%%%%%%%%%%%%%%%%%%%%%%%%%%%%%%%%%%

We say

\[
\boxed{
\lim_{z\to z_0}f(z)=L
}
\]

if for every \(\varepsilon>0\), there exists \(\delta>0\) such that

\[
0<|z-z_0|<\delta
\]

implies

\[
|f(z)-L|<\varepsilon.
\]

The symbols \(\varepsilon\) and \(\delta\) are lowercase Greek
\emph{epsilon} and \emph{delta}.

Because \(z\) moves in the plane, the limit must be the same along every
possible path approaching \(z_0\).

%%%%%%%%%%%%%%%%%%%%%%%%%%%%%%%%%%%%%%%%%%%%%%%%%%%%%%%%%%%%
\subsection{Continuity}
\label{subsec:complex-continuity}
%%%%%%%%%%%%%%%%%%%%%%%%%%%%%%%%%%%%%%%%%%%%%%%%%%%%%%%%%%%%

A function \(f\) is continuous at \(z_0\) if

\[
\boxed{
\lim_{z\to z_0}f(z)
=
f(z_0).
}
\]

If

\[
f(z)=u(x,y)+iv(x,y),
\]

then \(f\) is continuous precisely when both \(u\) and \(v\) are
continuous as real-valued functions of two real variables.

%%%%%%%%%%%%%%%%%%%%%%%%%%%%%%%%%%%%%%%%%%%%%%%%%%%%%%%%%%%%
\subsection{Complex Differentiability}
\label{subsec:complex-differentiability}
%%%%%%%%%%%%%%%%%%%%%%%%%%%%%%%%%%%%%%%%%%%%%%%%%%%%%%%%%%%%

\begin{definitionbox}{Complex Derivative}
The derivative of \(f\) at \(z_0\) is

\[
\boxed{
f'(z_0)
=
\lim_{z\to z_0}
\frac{f(z)-f(z_0)}
{z-z_0},
}
\]

provided the limit exists.
\end{definitionbox}

Equivalently,

\[
\boxed{
f'(z_0)
=
\lim_{h\to0}
\frac{f(z_0+h)-f(z_0)}{h}.
}
\]

Unlike the real case, \(h\) may approach zero from infinitely many
directions in the complex plane.

This makes complex differentiability much more restrictive than real
differentiability.

%%%%%%%%%%%%%%%%%%%%%%%%%%%%%%%%%%%%%%%%%%%%%%%%%%%%%%%%%%%%
\subsection{Holomorphic Functions}
\label{subsec:holomorphic-functions}
%%%%%%%%%%%%%%%%%%%%%%%%%%%%%%%%%%%%%%%%%%%%%%%%%%%%%%%%%%%%

\begin{definitionbox}{Holomorphic Function}
A function

\[
f:U\to\mathbb{C}
\]

defined on an open set \(U\subseteq\mathbb{C}\) is holomorphic on \(U\) if
it is complex differentiable at every point of \(U\).
\end{definitionbox}

A function holomorphic on all of \(\mathbb{C}\) is called entire.

Examples include

\[
\boxed{
e^z,
\qquad
\sin z,
\qquad
\cos z,
\qquad
\text{polynomials}.
}
\]

%%%%%%%%%%%%%%%%%%%%%%%%%%%%%%%%%%%%%%%%%%%%%%%%%%%%%%%%%%%%
\subsection{Cauchy--Riemann Equations}
\label{subsec:cauchy-riemann}
%%%%%%%%%%%%%%%%%%%%%%%%%%%%%%%%%%%%%%%%%%%%%%%%%%%%%%%%%%%%

Suppose

\[
f(z)=u(x,y)+iv(x,y).
\]

If \(f\) is complex differentiable and the relevant partial derivatives
exist, then

\[
\boxed{
u_x=v_y
}
\]

and

\[
\boxed{
u_y=-v_x.
}
\]

These are the Cauchy--Riemann equations.

They express the compatibility required between the real and imaginary
parts of a complex differentiable function.

%%%%%%%%%%%%%%%%%%%%%%%%%%%%%%%%%%%%%%%%%%%%%%%%%%%%%%%%%%%%
\subsection{Derivation of the Cauchy--Riemann Equations}
\label{subsec:derivation-cauchy-riemann}
%%%%%%%%%%%%%%%%%%%%%%%%%%%%%%%%%%%%%%%%%%%%%%%%%%%%%%%%%%%%

Approach \(z\) horizontally by taking

\[
h\in\R.
\]

Then

\[
f'(z)
=
u_x+iv_x.
\]

Approach vertically by taking

\[
h=ik,
\qquad
k\in\R.
\]

Then

\[
f'(z)
=
v_y-iu_y.
\]

Since the derivative must be independent of the direction of approach,

\[
u_x+iv_x
=
v_y-iu_y.
\]

Equating real and imaginary parts gives

\[
\boxed{
u_x=v_y,
\qquad
u_y=-v_x.
}
\]

%%%%%%%%%%%%%%%%%%%%%%%%%%%%%%%%%%%%%%%%%%%%%%%%%%%%%%%%%%%%
\subsection{Sufficient Conditions}
\label{subsec:cauchy-riemann-sufficient}
%%%%%%%%%%%%%%%%%%%%%%%%%%%%%%%%%%%%%%%%%%%%%%%%%%%%%%%%%%%%

The Cauchy--Riemann equations at a single point are not by themselves
always sufficient for complex differentiability.

A standard sufficient condition is the following.

\begin{theorem}
Suppose \(u_x,u_y,v_x,v_y\) exist and are continuous in a neighborhood of
\(z_0\).

If the Cauchy--Riemann equations hold at \(z_0\), then \(f\) is complex
differentiable at \(z_0\).
\end{theorem}

If these conditions hold throughout an open set, then \(f\) is
holomorphic there.

%%%%%%%%%%%%%%%%%%%%%%%%%%%%%%%%%%%%%%%%%%%%%%%%%%%%%%%%%%%%
\subsection{Worked Example: Testing Holomorphicity}
\label{subsec:worked-holomorphicity}
%%%%%%%%%%%%%%%%%%%%%%%%%%%%%%%%%%%%%%%%%%%%%%%%%%%%%%%%%%%%

\begin{workedexamplebox}{The Function \(f(z)=z^2\)}

Let

\[
z=x+iy.
\]

Then

\[
f(z)
=
(x+iy)^2
=
x^2-y^2+2ixy.
\]

Thus

\[
u(x,y)=x^2-y^2,
\qquad
v(x,y)=2xy.
\]

Compute

\[
u_x=2x,
\qquad
u_y=-2y,
\]

and

\[
v_x=2y,
\qquad
v_y=2x.
\]

Therefore

\[
u_x=v_y
\]

and

\[
u_y=-v_x.
\]

Hence \(f\) is holomorphic on \(\mathbb{C}\).

Its derivative is

\[
\begin{answerbox}
\[
\boxed{
f'(z)=2z.
}
\]
\end{answerbox}
\end{workedexamplebox}

%%%%%%%%%%%%%%%%%%%%%%%%%%%%%%%%%%%%%%%%%%%%%%%%%%%%%%%%%%%%
\subsection{A Non-Holomorphic Example}
\label{subsec:nonholomorphic-conjugate}
%%%%%%%%%%%%%%%%%%%%%%%%%%%%%%%%%%%%%%%%%%%%%%%%%%%%%%%%%%%%

Consider

\[
f(z)=\overline{z}.
\]

Then

\[
f(z)=x-iy,
\]

so

\[
u=x,
\qquad
v=-y.
\]

Hence

\[
u_x=1,
\qquad
v_y=-1.
\]

The Cauchy--Riemann equation

\[
u_x=v_y
\]

fails everywhere.

Therefore

\[
\boxed{
f(z)=\overline{z}
}
\]

is nowhere holomorphic.

%%%%%%%%%%%%%%%%%%%%%%%%%%%%%%%%%%%%%%%%%%%%%%%%%%%%%%%%%%%%
\subsection{Harmonic Functions}
\label{subsec:harmonic-functions}
%%%%%%%%%%%%%%%%%%%%%%%%%%%%%%%%%%%%%%%%%%%%%%%%%%%%%%%%%%%%

\begin{definitionbox}{Harmonic Function}
A twice continuously differentiable real-valued function \(u(x,y)\) is
harmonic if

\[
\boxed{
u_{xx}+u_{yy}=0.
}
\]
\end{definitionbox}

The operator

\[
\boxed{
\Delta
=
\frac{\partial^2}{\partial x^2}
+
\frac{\partial^2}{\partial y^2}
}
\]

is called the Laplacian.

The symbol \(\Delta\) is uppercase Greek \emph{Delta}.

Thus harmonicity is written

\[
\boxed{
\Delta u=0.
}
\]

%%%%%%%%%%%%%%%%%%%%%%%%%%%%%%%%%%%%%%%%%%%%%%%%%%%%%%%%%%%%
\subsection{Holomorphic Functions and Harmonicity}
\label{subsec:holomorphic-harmonic}
%%%%%%%%%%%%%%%%%%%%%%%%%%%%%%%%%%%%%%%%%%%%%%%%%%%%%%%%%%%%

Suppose

\[
f=u+iv
\]

is holomorphic and has sufficiently smooth second derivatives.

From

\[
u_x=v_y,
\qquad
u_y=-v_x,
\]

differentiate to obtain

\[
u_{xx}=v_{yx},
\]

and

\[
u_{yy}=-v_{xy}.
\]

Therefore

\[
\boxed{
u_{xx}+u_{yy}=0.
}
\]

Similarly,

\[
\boxed{
v_{xx}+v_{yy}=0.
}
\]

Thus the real and imaginary parts of a holomorphic function are harmonic.

%%%%%%%%%%%%%%%%%%%%%%%%%%%%%%%%%%%%%%%%%%%%%%%%%%%%%%%%%%%%
\subsection{Harmonic Conjugates}
\label{subsec:harmonic-conjugates}
%%%%%%%%%%%%%%%%%%%%%%%%%%%%%%%%%%%%%%%%%%%%%%%%%%%%%%%%%%%%

If

\[
f=u+iv
\]

is holomorphic, then \(v\) is called a harmonic conjugate of \(u\).

The Cauchy--Riemann equations can be used to recover one component from
the other.

Given \(u\),

\[
v_y=u_x,
\qquad
v_x=-u_y.
\]

Integrating these relations can produce \(v\), up to an additive constant,
on suitable domains.

%%%%%%%%%%%%%%%%%%%%%%%%%%%%%%%%%%%%%%%%%%%%%%%%%%%%%%%%%%%%
\subsection{The Complex Exponential}
\label{subsec:complex-exponential}
%%%%%%%%%%%%%%%%%%%%%%%%%%%%%%%%%%%%%%%%%%%%%%%%%%%%%%%%%%%%

For

\[
z=x+iy,
\]

define

\[
\boxed{
e^z
=
e^xe^{iy}
=
e^x(\cos y+i\sin y).
}
\]

The complex exponential is entire and satisfies

\[
\boxed{
\frac{d}{dz}e^z=e^z.
}
\]

It also satisfies

\[
\boxed{
e^{z+w}=e^ze^w.
}
\]

Unlike the real exponential, it is periodic:

\[
\boxed{
e^{z+2\pi i}=e^z.
}
\]

%%%%%%%%%%%%%%%%%%%%%%%%%%%%%%%%%%%%%%%%%%%%%%%%%%%%%%%%%%%%
\subsection{Complex Trigonometric Functions}
\label{subsec:complex-trigonometric-functions}
%%%%%%%%%%%%%%%%%%%%%%%%%%%%%%%%%%%%%%%%%%%%%%%%%%%%%%%%%%%%

Complex sine and cosine can be defined by

\[
\boxed{
\cos z
=
\frac{e^{iz}+e^{-iz}}{2}
}
\]

and

\[
\boxed{
\sin z
=
\frac{e^{iz}-e^{-iz}}{2i}.
}
\]

These functions are entire.

Their derivatives retain the familiar forms

\[
\boxed{
(\sin z)'=\cos z,
}
\]

\[
\boxed{
(\cos z)'=-\sin z.
}
\]

%%%%%%%%%%%%%%%%%%%%%%%%%%%%%%%%%%%%%%%%%%%%%%%%%%%%%%%%%%%%
\subsection{The Complex Logarithm}
\label{subsec:complex-logarithm}
%%%%%%%%%%%%%%%%%%%%%%%%%%%%%%%%%%%%%%%%%%%%%%%%%%%%%%%%%%%%

If

\[
z=re^{i\theta},
\]

then formally

\[
\boxed{
\log z
=
\ln r+i\theta.
}
\]

Because

\[
\theta
\]

is determined only modulo \(2\pi\), the complex logarithm is naturally
multivalued:

\[
\boxed{
\log z
=
\ln|z|
+
i(\arg z+2\pi k),
\qquad
k\in\Z.
}
\]

%%%%%%%%%%%%%%%%%%%%%%%%%%%%%%%%%%%%%%%%%%%%%%%%%%%%%%%%%%%%
\subsection{Principal Logarithm}
\label{subsec:principal-logarithm}
%%%%%%%%%%%%%%%%%%%%%%%%%%%%%%%%%%%%%%%%%%%%%%%%%%%%%%%%%%%%

Using a principal argument, one defines a single-valued branch such as

\[
\boxed{
\operatorname{Log}z
=
\ln|z|
+
i\operatorname{Arg}z.
}
\]

Under the convention

\[
-\pi<\operatorname{Arg}z<\pi,
\]

the principal logarithm is holomorphic on

\[
\boxed{
\mathbb{C}\setminus(-\infty,0].
}
\]

The removed ray is called a branch cut.

%%%%%%%%%%%%%%%%%%%%%%%%%%%%%%%%%%%%%%%%%%%%%%%%%%%%%%%%%%%%
\subsection{Branches and Branch Cuts}
\label{subsec:branches-branch-cuts}
%%%%%%%%%%%%%%%%%%%%%%%%%%%%%%%%%%%%%%%%%%%%%%%%%%%%%%%%%%%%

Multivalued expressions such as

\[
\log z,
\qquad
z^\alpha,
\qquad
\sqrt{z}
\]

cannot generally be made globally single-valued and holomorphic on
\(\mathbb{C}\setminus\{0\}\).

A branch is a consistent single-valued choice on a suitable domain.

A branch cut removes points so that such a choice becomes possible.

This phenomenon reflects the topology of the punctured complex plane.

%%%%%%%%%%%%%%%%%%%%%%%%%%%%%%%%%%%%%%%%%%%%%%%%%%%%%%%%%%%%
\subsection{Complex Powers}
\label{subsec:complex-powers}
%%%%%%%%%%%%%%%%%%%%%%%%%%%%%%%%%%%%%%%%%%%%%%%%%%%%%%%%%%%%

Using a chosen branch of the logarithm, define

\[
\boxed{
z^\alpha
=
e^{\alpha\operatorname{Log}z}.
}
\]

Different choices of logarithm may produce different values.

For example,

\[
\sqrt{z}
=
e^{\frac12\log z}
\]

is naturally two-valued before a branch is chosen.

%%%%%%%%%%%%%%%%%%%%%%%%%%%%%%%%%%%%%%%%%%%%%%%%%%%%%%%%%%%%
\subsection{Curves in the Complex Plane}
\label{subsec:complex-curves}
%%%%%%%%%%%%%%%%%%%%%%%%%%%%%%%%%%%%%%%%%%%%%%%%%%%%%%%%%%%%

A parameterized curve is a function

\[
\boxed{
\gamma:[a,b]\to\mathbb{C}.
}
\]

Write

\[
\gamma(t)
=
x(t)+iy(t).
\]

Its derivative is

\[
\boxed{
\gamma'(t)
=
x'(t)+iy'(t).
}
\]

The lowercase Greek letter \(\gamma\) is \emph{gamma} and commonly denotes
a curve or path.

%%%%%%%%%%%%%%%%%%%%%%%%%%%%%%%%%%%%%%%%%%%%%%%%%%%%%%%%%%%%
\subsection{Contour Integrals}
\label{subsec:contour-integrals}
%%%%%%%%%%%%%%%%%%%%%%%%%%%%%%%%%%%%%%%%%%%%%%%%%%%%%%%%%%%%

\begin{definitionbox}{Complex Contour Integral}
Let \(f\) be continuous on the image of a piecewise smooth curve
\(\gamma:[a,b]\to\mathbb{C}\).

Define

\[
\boxed{
\int_\gamma f(z)\,dz
=
\int_a^b
f(\gamma(t))\gamma'(t)\,dt.
}
\]
\end{definitionbox}

Thus a complex contour integral reduces to an ordinary real-parameter
integral.

%%%%%%%%%%%%%%%%%%%%%%%%%%%%%%%%%%%%%%%%%%%%%%%%%%%%%%%%%%%%
\subsection{Worked Example: Contour Integral}
\label{subsec:worked-contour-integral}
%%%%%%%%%%%%%%%%%%%%%%%%%%%%%%%%%%%%%%%%%%%%%%%%%%%%%%%%%%%%

\begin{workedexamplebox}{Integrating \(z\) Along a Line Segment}

Let \(\gamma\) be the line segment from \(0\) to \(1+i\):

\[
\gamma(t)=t(1+i),
\qquad
0\leq t\leq1.
\]

Then

\[
\gamma'(t)=1+i.
\]

Therefore

\[
\int_\gamma z\,dz
=
\int_0^1
t(1+i)(1+i)\,dt.
\]

Since

\[
(1+i)^2=2i,
\]

we obtain

\[
\int_0^1 2it\,dt=i.
\]

Hence

\[
\begin{answerbox}
\[
\boxed{
\int_\gamma z\,dz=i.
}
\]
\end{answerbox}
\end{workedexamplebox}

%%%%%%%%%%%%%%%%%%%%%%%%%%%%%%%%%%%%%%%%%%%%%%%%%%%%%%%%%%%%
\subsection{Basic Properties of Contour Integrals}
\label{subsec:properties-contour-integrals}
%%%%%%%%%%%%%%%%%%%%%%%%%%%%%%%%%%%%%%%%%%%%%%%%%%%%%%%%%%%%

Contour integrals are linear:

\[
\boxed{
\int_\gamma(af+bg)\,dz
=
a\int_\gamma f\,dz
+
b\int_\gamma g\,dz.
}
\]

Reversing orientation changes the sign:

\[
\boxed{
\int_{-\gamma}f(z)\,dz
=
-\int_\gamma f(z)\,dz.
}
\]

If a curve is formed by concatenation,

\[
\gamma=\gamma_1*\gamma_2,
\]

then

\[
\boxed{
\int_\gamma f\,dz
=
\int_{\gamma_1}f\,dz
+
\int_{\gamma_2}f\,dz.
}
\]

%%%%%%%%%%%%%%%%%%%%%%%%%%%%%%%%%%%%%%%%%%%%%%%%%%%%%%%%%%%%
\subsection{The ML Estimate}
\label{subsec:ml-estimate}
%%%%%%%%%%%%%%%%%%%%%%%%%%%%%%%%%%%%%%%%%%%%%%%%%%%%%%%%%%%%

Suppose

\[
|f(z)|\leq M
\]

along a contour \(\gamma\) of length \(L\).

Then

\[
\boxed{
\left|
\int_\gamma f(z)\,dz
\right|
\leq
ML.
}
\]

This is commonly called the ML estimate or estimation lemma.

It is extremely useful for proving that contour integrals tend to zero.

%%%%%%%%%%%%%%%%%%%%%%%%%%%%%%%%%%%%%%%%%%%%%%%%%%%%%%%%%%%%
\subsection{Antiderivatives and Path Independence}
\label{subsec:complex-antiderivatives}
%%%%%%%%%%%%%%%%%%%%%%%%%%%%%%%%%%%%%%%%%%%%%%%%%%%%%%%%%%%%

If

\[
F'(z)=f(z),
\]

then \(F\) is an antiderivative of \(f\).

For a path from \(z_0\) to \(z_1\),

\[
\boxed{
\int_\gamma f(z)\,dz
=
F(z_1)-F(z_0).
}
\]

Consequently, the integral depends only on the endpoints and not on the
specific path.

In particular, for every closed contour,

\[
\boxed{
\oint_\gamma f(z)\,dz=0.
}
\]

The symbol

\[
\oint
\]

denotes integration around a closed contour.

%%%%%%%%%%%%%%%%%%%%%%%%%%%%%%%%%%%%%%%%%%%%%%%%%%%%%%%%%%%%
\subsection{Cauchy's Integral Theorem}
\label{subsec:cauchy-integral-theorem}
%%%%%%%%%%%%%%%%%%%%%%%%%%%%%%%%%%%%%%%%%%%%%%%%%%%%%%%%%%%%

One of the central results of complex analysis is that holomorphic
functions have vanishing integrals around suitable closed curves.

\begin{theorem}[Cauchy's Integral Theorem]
Let \(D\subseteq\mathbb{C}\) be a simply connected domain and let
\(f:D\to\mathbb{C}\) be holomorphic.

Then for every closed piecewise smooth contour \(\gamma\) in \(D\),

\[
\boxed{
\oint_\gamma f(z)\,dz=0.
}
\]
\end{theorem}

This result has no direct analogue of comparable strength in elementary
real calculus.

%%%%%%%%%%%%%%%%%%%%%%%%%%%%%%%%%%%%%%%%%%%%%%%%%%%%%%%%%%%%
\subsection{Simply Connected Domains}
\label{subsec:simply-connected-complex}
%%%%%%%%%%%%%%%%%%%%%%%%%%%%%%%%%%%%%%%%%%%%%%%%%%%%%%%%%%%%

Informally, a domain is simply connected if it has no holes.

Every closed curve in the domain can be continuously contracted to a
point while remaining inside the domain.

For example,

\[
\mathbb{C}
\]

is simply connected.

However,

\[
\mathbb{C}\setminus\{0\}
\]

is not simply connected.

This distinction is crucial for functions such as

\[
\frac1z.
\]

%%%%%%%%%%%%%%%%%%%%%%%%%%%%%%%%%%%%%%%%%%%%%%%%%%%%%%%%%%%%
\subsection{Integral of \(1/z\)}
\label{subsec:integral-one-over-z}
%%%%%%%%%%%%%%%%%%%%%%%%%%%%%%%%%%%%%%%%%%%%%%%%%%%%%%%%%%%%

Consider the positively oriented circle

\[
\gamma(t)=Re^{it},
\qquad
0\leq t\leq2\pi.
\]

Then

\[
\gamma'(t)=iRe^{it}.
\]

Therefore

\[
\oint_\gamma\frac1z\,dz
=
\int_0^{2\pi}
\frac{1}{Re^{it}}
iRe^{it}\,dt.
\]

Hence

\[
\boxed{
\oint_\gamma\frac1z\,dz
=
2\pi i.
}
\]

The integral is nonzero because \(1/z\) has a singularity at the origin,
which lies inside the contour.

%%%%%%%%%%%%%%%%%%%%%%%%%%%%%%%%%%%%%%%%%%%%%%%%%%%%%%%%%%%%
\subsection{Cauchy's Integral Formula}
\label{subsec:cauchy-integral-formula}
%%%%%%%%%%%%%%%%%%%%%%%%%%%%%%%%%%%%%%%%%%%%%%%%%%%%%%%%%%%%

\begin{theorem}[Cauchy's Integral Formula]
Suppose \(f\) is holomorphic on a domain containing a positively oriented
simple closed contour \(\gamma\) and its interior.

If \(a\) lies inside \(\gamma\), then

\[
\boxed{
f(a)
=
\frac{1}{2\pi i}
\oint_\gamma
\frac{f(z)}{z-a}\,dz.
}
\]
\end{theorem}

Equivalently,

\[
\boxed{
\oint_\gamma
\frac{f(z)}{z-a}\,dz
=
2\pi i\,f(a).
}
\]

This remarkable formula recovers the value of a holomorphic function
inside a contour entirely from its boundary values.

%%%%%%%%%%%%%%%%%%%%%%%%%%%%%%%%%%%%%%%%%%%%%%%%%%%%%%%%%%%%
\subsection{Worked Example: Cauchy's Integral Formula}
\label{subsec:worked-cauchy-formula}
%%%%%%%%%%%%%%%%%%%%%%%%%%%%%%%%%%%%%%%%%%%%%%%%%%%%%%%%%%%%

\begin{workedexamplebox}{Evaluating a Contour Integral}

Evaluate

\[
\oint_{|z|=2}
\frac{e^z}{z-1}\,dz
\]

with positive orientation.

The point

\[
a=1
\]

lies inside the circle.

Take

\[
f(z)=e^z.
\]

By Cauchy's Integral Formula,

\[
\oint_{|z|=2}
\frac{e^z}{z-1}\,dz
=
2\pi i\,e.
\]

Therefore

\[
\begin{answerbox}
\[
\boxed{
\oint_{|z|=2}
\frac{e^z}{z-1}\,dz
=
2\pi i e.
}
\]
\end{answerbox}
\end{workedexamplebox}

%%%%%%%%%%%%%%%%%%%%%%%%%%%%%%%%%%%%%%%%%%%%%%%%%%%%%%%%%%%%
\subsection{Cauchy's Formula for Derivatives}
\label{subsec:cauchy-derivative-formula}
%%%%%%%%%%%%%%%%%%%%%%%%%%%%%%%%%%%%%%%%%%%%%%%%%%%%%%%%%%%%

Cauchy's Integral Formula extends to derivatives.

\begin{theorem}
Under the hypotheses of Cauchy's Integral Formula,

\[
\boxed{
f^{(n)}(a)
=
\frac{n!}{2\pi i}
\oint_\gamma
\frac{f(z)}
{(z-a)^{n+1}}
\,dz.
}
\]
\end{theorem}

Thus a holomorphic function automatically has derivatives of every order.

This is a dramatic contrast with real analysis.

%%%%%%%%%%%%%%%%%%%%%%%%%%%%%%%%%%%%%%%%%%%%%%%%%%%%%%%%%%%%
\subsection{Holomorphic Implies Analytic}
\label{subsec:holomorphic-implies-analytic}
%%%%%%%%%%%%%%%%%%%%%%%%%%%%%%%%%%%%%%%%%%%%%%%%%%%%%%%%%%%%

A function is analytic near \(a\) if it can be represented by a convergent
power series

\[
\boxed{
f(z)
=
\sum_{n=0}^{\infty}
a_n(z-a)^n.
}
\]

A fundamental theorem of complex analysis states

\[
\boxed{
\text{holomorphic}
\Longrightarrow
\text{analytic}.
}
\]

In fact, the coefficients are

\[
\boxed{
a_n
=
\frac{f^{(n)}(a)}{n!}.
}
\]

Therefore

\[
\boxed{
f(z)
=
\sum_{n=0}^{\infty}
\frac{f^{(n)}(a)}{n!}
(z-a)^n.
}
\]

%%%%%%%%%%%%%%%%%%%%%%%%%%%%%%%%%%%%%%%%%%%%%%%%%%%%%%%%%%%%
\subsection{Cauchy Estimates}
\label{subsec:cauchy-estimates}
%%%%%%%%%%%%%%%%%%%%%%%%%%%%%%%%%%%%%%%%%%%%%%%%%%%%%%%%%%%%

Suppose \(f\) is holomorphic on and inside

\[
|z-a|=R
\]

and

\[
|f(z)|\leq M
\]

on the circle.

Then Cauchy's derivative formula yields

\[
\boxed{
|f^{(n)}(a)|
\leq
\frac{n!M}{R^n}.
}
\]

These inequalities are called Cauchy estimates.

They are useful for controlling derivatives of holomorphic functions.

%%%%%%%%%%%%%%%%%%%%%%%%%%%%%%%%%%%%%%%%%%%%%%%%%%%%%%%%%%%%
\subsection{Liouville's Theorem}
\label{subsec:liouvilles-theorem}
%%%%%%%%%%%%%%%%%%%%%%%%%%%%%%%%%%%%%%%%%%%%%%%%%%%%%%%%%%%%

\begin{theorem}[Liouville's Theorem]
Every bounded entire function is constant.
\end{theorem}

That is, if \(f\) is holomorphic on all of \(\mathbb{C}\) and there exists
\(M>0\) such that

\[
|f(z)|\leq M
\]

for every \(z\), then

\[
\boxed{
f(z)=C
}
\]

for some constant \(C\).

The proof follows immediately from Cauchy estimates by letting the radius
of the bounding circle tend to infinity.

%%%%%%%%%%%%%%%%%%%%%%%%%%%%%%%%%%%%%%%%%%%%%%%%%%%%%%%%%%%%
\subsection{Fundamental Theorem of Algebra}
\label{subsec:fundamental-theorem-algebra-complex}
%%%%%%%%%%%%%%%%%%%%%%%%%%%%%%%%%%%%%%%%%%%%%%%%%%%%%%%%%%%%

Liouville's Theorem gives an elegant proof of a central algebraic result.

\begin{theorem}[Fundamental Theorem of Algebra]
Every nonconstant polynomial with complex coefficients has at least one
complex root.
\end{theorem}

Consequently, every degree-\(n\) polynomial over \(\mathbb{C}\) factors as

\[
\boxed{
p(z)
=
a
\prod_{k=1}^{n}(z-z_k),
}
\]

counting roots with multiplicity.

Thus

\[
\boxed{
\mathbb{C}
\text{ is algebraically closed}.
}
\]

%%%%%%%%%%%%%%%%%%%%%%%%%%%%%%%%%%%%%%%%%%%%%%%%%%%%%%%%%%%%
\subsection{Maximum Modulus Principle}
\label{subsec:maximum-modulus-principle}
%%%%%%%%%%%%%%%%%%%%%%%%%%%%%%%%%%%%%%%%%%%%%%%%%%%%%%%%%%%%

\begin{theorem}[Maximum Modulus Principle]
Let \(f\) be a nonconstant holomorphic function on a connected open set
\(D\).

Then \(|f|\) cannot attain a local maximum at an interior point of \(D\).
\end{theorem}

For a bounded domain with suitable boundary behavior, this means the
maximum of \(|f|\) occurs on the boundary.

This principle is one of the strongest rigidity results in complex
analysis.

%%%%%%%%%%%%%%%%%%%%%%%%%%%%%%%%%%%%%%%%%%%%%%%%%%%%%%%%%%%%
\subsection{Identity Theorem}
\label{subsec:identity-theorem-complex}
%%%%%%%%%%%%%%%%%%%%%%%%%%%%%%%%%%%%%%%%%%%%%%%%%%%%%%%%%%%%

\begin{theorem}[Identity Theorem]
Let \(f\) and \(g\) be holomorphic on a connected domain \(D\).

If

\[
f(z)=g(z)
\]

on a set having an accumulation point inside \(D\), then

\[
\boxed{
f\equiv g
}
\]

throughout \(D\).
\end{theorem}

Thus holomorphic functions are completely determined by surprisingly
small sets of values.

%%%%%%%%%%%%%%%%%%%%%%%%%%%%%%%%%%%%%%%%%%%%%%%%%%%%%%%%%%%%
\subsection{Zeros of Holomorphic Functions}
\label{subsec:zeros-holomorphic}
%%%%%%%%%%%%%%%%%%%%%%%%%%%%%%%%%%%%%%%%%%%%%%%%%%%%%%%%%%%%

Suppose \(f\) is holomorphic near \(a\) and

\[
f(a)=0.
\]

If \(f\) is not identically zero, then there exists an integer \(m\geq1\)
and a holomorphic function \(g\) with

\[
g(a)\neq0
\]

such that

\[
\boxed{
f(z)
=
(z-a)^m g(z).
}
\]

The integer \(m\) is the multiplicity or order of the zero.

Zeros of a nonzero holomorphic function are isolated.

%%%%%%%%%%%%%%%%%%%%%%%%%%%%%%%%%%%%%%%%%%%%%%%%%%%%%%%%%%%%
\subsection{Laurent Series}
\label{subsec:laurent-series}
%%%%%%%%%%%%%%%%%%%%%%%%%%%%%%%%%%%%%%%%%%%%%%%%%%%%%%%%%%%%

Taylor series use nonnegative powers.

Near singularities, negative powers may also be necessary.

\begin{definitionbox}{Laurent Series}
A Laurent series centered at \(a\) has the form

\[
\boxed{
\sum_{n=-\infty}^{\infty}
c_n(z-a)^n.
}
\]
\end{definitionbox}

It may be written as

\[
\boxed{
\sum_{n=0}^{\infty}
c_n(z-a)^n
+
\sum_{n=1}^{\infty}
\frac{c_{-n}}{(z-a)^n}.
}
\]

The second sum is called the principal part.

%%%%%%%%%%%%%%%%%%%%%%%%%%%%%%%%%%%%%%%%%%%%%%%%%%%%%%%%%%%%
\subsection{Laurent Coefficients}
\label{subsec:laurent-coefficients}
%%%%%%%%%%%%%%%%%%%%%%%%%%%%%%%%%%%%%%%%%%%%%%%%%%%%%%%%%%%%

If \(f\) is holomorphic on an annulus containing a positively oriented
circle \(\gamma\) around \(a\), then

\[
\boxed{
c_n
=
\frac{1}{2\pi i}
\oint_\gamma
\frac{f(z)}
{(z-a)^{n+1}}
\,dz.
}
\]

Thus Laurent coefficients generalize Taylor coefficients.

%%%%%%%%%%%%%%%%%%%%%%%%%%%%%%%%%%%%%%%%%%%%%%%%%%%%%%%%%%%%
\subsection{Isolated Singularities}
\label{subsec:isolated-singularities}
%%%%%%%%%%%%%%%%%%%%%%%%%%%%%%%%%%%%%%%%%%%%%%%%%%%%%%%%%%%%

Suppose \(f\) is holomorphic in a punctured neighborhood

\[
0<|z-a|<R.
\]

Then \(a\) is called an isolated singularity.

There are three principal types:

\[
\boxed{
\text{removable singularity},
}
\]

\[
\boxed{
\text{pole},
}
\]

and

\[
\boxed{
\text{essential singularity}.
}
\]

The Laurent series determines the type.

%%%%%%%%%%%%%%%%%%%%%%%%%%%%%%%%%%%%%%%%%%%%%%%%%%%%%%%%%%%%
\subsection{Removable Singularities}
\label{subsec:removable-singularities}
%%%%%%%%%%%%%%%%%%%%%%%%%%%%%%%%%%%%%%%%%%%%%%%%%%%%%%%%%%%%

A singularity at \(a\) is removable if the principal part of the Laurent
series vanishes.

Equivalently,

\[
\boxed{
c_{-1}=c_{-2}=\cdots=0.
}
\]

The function can then be extended holomorphically to \(a\).

For example,

\[
f(z)
=
\frac{\sin z}{z}
\]

has a removable singularity at \(z=0\), since

\[
\lim_{z\to0}\frac{\sin z}{z}=1.
\]

Defining

\[
f(0)=1
\]

removes the singularity.

%%%%%%%%%%%%%%%%%%%%%%%%%%%%%%%%%%%%%%%%%%%%%%%%%%%%%%%%%%%%
\subsection{Poles}
\label{subsec:complex-poles}
%%%%%%%%%%%%%%%%%%%%%%%%%%%%%%%%%%%%%%%%%%%%%%%%%%%%%%%%%%%%

A singularity \(a\) is a pole of order \(m\) if the Laurent series has
finitely many negative powers and the highest is

\[
\frac{c_{-m}}{(z-a)^m},
\qquad
c_{-m}\neq0.
\]

Equivalently,

\[
\boxed{
f(z)
=
\frac{g(z)}
{(z-a)^m},
}
\]

where \(g\) is holomorphic near \(a\) and

\[
g(a)\neq0.
\]

%%%%%%%%%%%%%%%%%%%%%%%%%%%%%%%%%%%%%%%%%%%%%%%%%%%%%%%%%%%%
\subsection{Essential Singularities}
\label{subsec:essential-singularities}
%%%%%%%%%%%%%%%%%%%%%%%%%%%%%%%%%%%%%%%%%%%%%%%%%%%%%%%%%%%%

A singularity is essential if its Laurent series contains infinitely many
negative-power terms.

For example,

\[
\boxed{
e^{1/z}
=
1+\frac1z+\frac1{2!z^2}
+\frac1{3!z^3}
+\cdots
}
\]

has an essential singularity at \(z=0\).

Essential singularities exhibit extremely complicated local behavior.

%%%%%%%%%%%%%%%%%%%%%%%%%%%%%%%%%%%%%%%%%%%%%%%%%%%%%%%%%%%%
\subsection{Residues}
\label{subsec:residues}
%%%%%%%%%%%%%%%%%%%%%%%%%%%%%%%%%%%%%%%%%%%%%%%%%%%%%%%%%%%%

\begin{definitionbox}{Residue}
If \(f\) has Laurent expansion

\[
f(z)
=
\sum_{n=-\infty}^{\infty}
c_n(z-a)^n,
\]

then the residue of \(f\) at \(a\) is

\[
\boxed{
\operatorname{Res}(f;a)
=
c_{-1}.
}
\]
\end{definitionbox}

Thus the residue is the coefficient of

\[
\frac1{z-a}.
\]

Although it is only one Laurent coefficient, it determines the contour
integral around the singularity.

%%%%%%%%%%%%%%%%%%%%%%%%%%%%%%%%%%%%%%%%%%%%%%%%%%%%%%%%%%%%
\subsection{Residue at a Simple Pole}
\label{subsec:residue-simple-pole}
%%%%%%%%%%%%%%%%%%%%%%%%%%%%%%%%%%%%%%%%%%%%%%%%%%%%%%%%%%%%

If \(a\) is a simple pole, then

\[
\boxed{
\operatorname{Res}(f;a)
=
\lim_{z\to a}
(z-a)f(z).
}
\]

If

\[
f(z)=\frac{g(z)}{h(z)},
\]

where

\[
h(a)=0,
\qquad
h'(a)\neq0,
\]

then

\[
\boxed{
\operatorname{Res}(f;a)
=
\frac{g(a)}{h'(a)}.
}
\]

%%%%%%%%%%%%%%%%%%%%%%%%%%%%%%%%%%%%%%%%%%%%%%%%%%%%%%%%%%%%
\subsection{Residue at a Higher-Order Pole}
\label{subsec:residue-higher-order}
%%%%%%%%%%%%%%%%%%%%%%%%%%%%%%%%%%%%%%%%%%%%%%%%%%%%%%%%%%%%

If \(a\) is a pole of order \(m\), then

\[
\boxed{
\operatorname{Res}(f;a)
=
\frac{1}{(m-1)!}
\lim_{z\to a}
\frac{d^{m-1}}{dz^{m-1}}
\left[
(z-a)^m f(z)
\right].
}
\]

This formula avoids constructing the full Laurent series.

%%%%%%%%%%%%%%%%%%%%%%%%%%%%%%%%%%%%%%%%%%%%%%%%%%%%%%%%%%%%
\subsection{Residue Theorem}
\label{subsec:residue-theorem}
%%%%%%%%%%%%%%%%%%%%%%%%%%%%%%%%%%%%%%%%%%%%%%%%%%%%%%%%%%%%

\begin{theorem}[Residue Theorem]
Let \(f\) be holomorphic inside and on a positively oriented simple closed
contour \(\gamma\), except for finitely many isolated singularities

\[
a_1,\ldots,a_n
\]

inside \(\gamma\).

Then

\[
\boxed{
\oint_\gamma f(z)\,dz
=
2\pi i
\sum_{k=1}^{n}
\operatorname{Res}(f;a_k).
}
\]
\end{theorem}

This theorem converts a contour integral into a finite algebraic
calculation.

%%%%%%%%%%%%%%%%%%%%%%%%%%%%%%%%%%%%%%%%%%%%%%%%%%%%%%%%%%%%
\subsection{Worked Example: Residue Theorem}
\label{subsec:worked-residue-theorem}
%%%%%%%%%%%%%%%%%%%%%%%%%%%%%%%%%%%%%%%%%%%%%%%%%%%%%%%%%%%%

\begin{workedexamplebox}{Contour Integral by Residues}

Evaluate

\[
\oint_{|z|=2}
\frac{1}{z(z-1)}
\,dz.
\]

The singularities inside the contour are

\[
z=0
\]

and

\[
z=1.
\]

At \(z=0\),

\[
\operatorname{Res}(f;0)
=
\lim_{z\to0}
\frac{1}{z-1}
=
-1.
\]

At \(z=1\),

\[
\operatorname{Res}(f;1)
=
\lim_{z\to1}
\frac{1}{z}
=
1.
\]

Therefore

\[
\sum\operatorname{Res}(f;a_k)
=
0.
\]

By the Residue Theorem,

\[
\begin{answerbox}
\[
\boxed{
\oint_{|z|=2}
\frac{1}{z(z-1)}
\,dz
=
0.
}
\]
\end{answerbox}
\end{workedexamplebox}

%%%%%%%%%%%%%%%%%%%%%%%%%%%%%%%%%%%%%%%%%%%%%%%%%%%%%%%%%%%%
\subsection{Real Integrals Using Residues}
\label{subsec:real-integrals-residues}
%%%%%%%%%%%%%%%%%%%%%%%%%%%%%%%%%%%%%%%%%%%%%%%%%%%%%%%%%%%%

Residues can evaluate certain improper real integrals.

A standard example is

\[
\int_{-\infty}^{\infty}
\frac{dx}{x^2+1}.
\]

Consider

\[
f(z)
=
\frac{1}{z^2+1}
=
\frac{1}{(z-i)(z+i)}.
\]

Using a semicircular contour in the upper half-plane, only the pole

\[
z=i
\]

lies inside.

Its residue is

\[
\operatorname{Res}(f;i)
=
\frac{1}{2i}.
\]

Hence

\[
2\pi i
\operatorname{Res}(f;i)
=
\pi.
\]

The semicircular contribution tends to zero as the radius tends to
infinity, giving

\[
\boxed{
\int_{-\infty}^{\infty}
\frac{dx}{x^2+1}
=
\pi.
}
\]

%%%%%%%%%%%%%%%%%%%%%%%%%%%%%%%%%%%%%%%%%%%%%%%%%%%%%%%%%%%%
\subsection{Meromorphic Functions}
\label{subsec:meromorphic-functions}
%%%%%%%%%%%%%%%%%%%%%%%%%%%%%%%%%%%%%%%%%%%%%%%%%%%%%%%%%%%%

\begin{definitionbox}{Meromorphic Function}
A function is meromorphic on a domain if it is holomorphic except at
isolated poles.
\end{definitionbox}

Rational functions are meromorphic on \(\mathbb{C}\).

For example,

\[
\boxed{
f(z)=\frac{1}{z^2+1}
}
\]

is meromorphic with poles at

\[
z=i,
\qquad
z=-i.
\]

%%%%%%%%%%%%%%%%%%%%%%%%%%%%%%%%%%%%%%%%%%%%%%%%%%%%%%%%%%%%
\subsection{Argument Principle}
\label{subsec:argument-principle}
%%%%%%%%%%%%%%%%%%%%%%%%%%%%%%%%%%%%%%%%%%%%%%%%%%%%%%%%%%%%

\begin{theorem}[Argument Principle]
Suppose \(f\) is meromorphic inside and on a positively oriented contour
\(\gamma\), with no zeros or poles on \(\gamma\).

Then

\[
\boxed{
\frac{1}{2\pi i}
\oint_\gamma
\frac{f'(z)}{f(z)}
\,dz
=
N-P,
}
\]

where \(N\) is the number of zeros and \(P\) is the number of poles inside
\(\gamma\), counted with multiplicity.
\end{theorem}

The theorem relates complex integration to the winding behavior of the
image curve \(f(\gamma)\).

%%%%%%%%%%%%%%%%%%%%%%%%%%%%%%%%%%%%%%%%%%%%%%%%%%%%%%%%%%%%
\subsection{Winding Number}
\label{subsec:winding-number-complex}
%%%%%%%%%%%%%%%%%%%%%%%%%%%%%%%%%%%%%%%%%%%%%%%%%%%%%%%%%%%%

The winding number of a closed curve \(\gamma\) about a point \(a\) not on
the curve is

\[
\boxed{
\operatorname{Ind}(\gamma,a)
=
\frac{1}{2\pi i}
\oint_\gamma
\frac{dz}{z-a}.
}
\]

It counts, with orientation, how many times the curve winds around \(a\).

For a positively oriented simple circle containing \(a\),

\[
\boxed{
\operatorname{Ind}(\gamma,a)=1.
}
\]

%%%%%%%%%%%%%%%%%%%%%%%%%%%%%%%%%%%%%%%%%%%%%%%%%%%%%%%%%%%%
\subsection{Rouch\'e's Theorem}
\label{subsec:rouches-theorem}
%%%%%%%%%%%%%%%%%%%%%%%%%%%%%%%%%%%%%%%%%%%%%%%%%%%%%%%%%%%%

\begin{theorem}[Rouch\'e's Theorem]
Let \(f\) and \(g\) be holomorphic inside and on a simple closed contour
\(\gamma\).

If

\[
\boxed{
|g(z)|<|f(z)|
}
\]

for every \(z\in\gamma\), then \(f\) and \(f+g\) have the same number of
zeros inside \(\gamma\), counted with multiplicity.
\end{theorem}

Rouch\'e's Theorem is a powerful method for counting zeros without solving
an equation explicitly.

%%%%%%%%%%%%%%%%%%%%%%%%%%%%%%%%%%%%%%%%%%%%%%%%%%%%%%%%%%%%
\subsection{Open Mapping Theorem}
\label{subsec:open-mapping-theorem-complex}
%%%%%%%%%%%%%%%%%%%%%%%%%%%%%%%%%%%%%%%%%%%%%%%%%%%%%%%%%%%%

\begin{theorem}[Open Mapping Theorem]
A nonconstant holomorphic function maps open sets to open sets.
\end{theorem}

Thus if \(D\subseteq\mathbb{C}\) is open and \(f\) is nonconstant and
holomorphic on \(D\), then

\[
\boxed{
f(D)
\text{ is open}.
}
\]

This is another strong rigidity property of holomorphic functions.

%%%%%%%%%%%%%%%%%%%%%%%%%%%%%%%%%%%%%%%%%%%%%%%%%%%%%%%%%%%%
\subsection{Conformal Maps}
\label{subsec:conformal-maps}
%%%%%%%%%%%%%%%%%%%%%%%%%%%%%%%%%%%%%%%%%%%%%%%%%%%%%%%%%%%%

A conformal map preserves oriented angles locally.

If \(f\) is holomorphic at \(z_0\) and

\[
\boxed{
f'(z_0)\neq0,
}
\]

then \(f\) is conformal at \(z_0\).

Locally, multiplication by

\[
f'(z_0)
=
re^{i\theta}
\]

acts as a scaling by \(r\) and a rotation by \(\theta\).

Thus complex differentiability has a direct geometric interpretation.

%%%%%%%%%%%%%%%%%%%%%%%%%%%%%%%%%%%%%%%%%%%%%%%%%%%%%%%%%%%%
\subsection{M\"obius Transformations}
\label{subsec:mobius-transformations}
%%%%%%%%%%%%%%%%%%%%%%%%%%%%%%%%%%%%%%%%%%%%%%%%%%%%%%%%%%%%

\begin{definitionbox}{M\"obius Transformation}
A M\"obius transformation has the form

\[
\boxed{
T(z)
=
\frac{az+b}{cz+d},
}
\]

where

\[
a,b,c,d\in\mathbb{C}
\]

and

\[
\boxed{
ad-bc\neq0.
}
\]
\end{definitionbox}

These transformations naturally act on the extended complex plane

\[
\boxed{
\widehat{\mathbb{C}}
=
\mathbb{C}\cup\{\infty\}.
}
\]

They map generalized circles---circles and lines---to generalized circles.

%%%%%%%%%%%%%%%%%%%%%%%%%%%%%%%%%%%%%%%%%%%%%%%%%%%%%%%%%%%%
\subsection{The Riemann Sphere}
\label{subsec:riemann-sphere}
%%%%%%%%%%%%%%%%%%%%%%%%%%%%%%%%%%%%%%%%%%%%%%%%%%%%%%%%%%%%

The extended complex plane

\[
\widehat{\mathbb{C}}
=
\mathbb{C}\cup\{\infty\}
\]

can be identified with a sphere using stereographic projection.

This geometric model is called the Riemann sphere.

It allows infinity to be treated as a single geometric point.

Rational functions then become natural maps

\[
\boxed{
\widehat{\mathbb{C}}
\to
\widehat{\mathbb{C}}.
}
\]

%%%%%%%%%%%%%%%%%%%%%%%%%%%%%%%%%%%%%%%%%%%%%%%%%%%%%%%%%%%%
\subsection{Schwarz Lemma}
\label{subsec:schwarz-lemma}
%%%%%%%%%%%%%%%%%%%%%%%%%%%%%%%%%%%%%%%%%%%%%%%%%%%%%%%%%%%%

Let

\[
\mathbb{D}
=
\{z\in\mathbb{C}:|z|<1\}
\]

denote the open unit disk.

\begin{theorem}[Schwarz Lemma]
Suppose

\[
f:\mathbb{D}\to\mathbb{D}
\]

is holomorphic and

\[
f(0)=0.
\]

Then

\[
\boxed{
|f(z)|\leq|z|
}
\]

for all \(z\in\mathbb{D}\), and

\[
\boxed{
|f'(0)|\leq1.
}
\]

If equality occurs at a nonzero point or in the derivative condition, then

\[
f(z)=e^{i\theta}z
\]

for some real \(\theta\).
\end{theorem}

%%%%%%%%%%%%%%%%%%%%%%%%%%%%%%%%%%%%%%%%%%%%%%%%%%%%%%%%%%%%
\subsection{Analytic Continuation}
\label{subsec:analytic-continuation}
%%%%%%%%%%%%%%%%%%%%%%%%%%%%%%%%%%%%%%%%%%%%%%%%%%%%%%%%%%%%

Suppose a holomorphic function is initially defined on a domain \(D_1\).

If there exists a holomorphic function on a larger domain \(D_2\) agreeing
with the original function on \(D_1\), then the larger function is an
analytic continuation.

The Identity Theorem implies that analytic continuation, when it exists,
is highly rigid.

This idea becomes central in advanced complex analysis and special
functions.

%%%%%%%%%%%%%%%%%%%%%%%%%%%%%%%%%%%%%%%%%%%%%%%%%%%%%%%%%%%%
\subsection{Common Errors}
\label{subsec:complex-analysis-common-errors}
%%%%%%%%%%%%%%%%%%%%%%%%%%%%%%%%%%%%%%%%%%%%%%%%%%%%%%%%%%%%

\subsubsection{Treating Complex Differentiability Like Real Differentiability}

The difference quotient must have the same limit for every direction of
approach in the complex plane.

%%%%%%%%%%%%%%%%%%%%%%%%%%%%%%%%%%%%%%%%%%%%%%%%%%%%%%%%%%%%
\subsubsection{Using the Cauchy--Riemann Equations Without Checking Hypotheses}

Satisfying the Cauchy--Riemann equations at one point does not
automatically guarantee complex differentiability there.

%%%%%%%%%%%%%%%%%%%%%%%%%%%%%%%%%%%%%%%%%%%%%%%%%%%%%%%%%%%%
\subsubsection{Treating the Argument as Single-Valued}

In general,

\[
\arg z
=
\theta+2\pi k.
\]

A principal argument is a selected branch.

%%%%%%%%%%%%%%%%%%%%%%%%%%%%%%%%%%%%%%%%%%%%%%%%%%%%%%%%%%%%
\subsubsection{Forgetting Branch Cuts}

Expressions such as

\[
\log z
\]

and

\[
z^\alpha
\]

require branch choices when treated as single-valued holomorphic
functions.

%%%%%%%%%%%%%%%%%%%%%%%%%%%%%%%%%%%%%%%%%%%%%%%%%%%%%%%%%%%%
\subsubsection{Applying Cauchy's Theorem Across a Singularity}

The function must be holomorphic throughout the relevant region.

For example,

\[
\frac1z
\]

is not holomorphic at the origin.

%%%%%%%%%%%%%%%%%%%%%%%%%%%%%%%%%%%%%%%%%%%%%%%%%%%%%%%%%%%%
\subsubsection{Ignoring Contour Orientation}

Positive orientation is conventionally counterclockwise.

Reversing orientation multiplies the integral by \(-1\).

%%%%%%%%%%%%%%%%%%%%%%%%%%%%%%%%%%%%%%%%%%%%%%%%%%%%%%%%%%%%
\subsubsection{Using Cauchy's Integral Formula When the Point Is Outside}

The point \(a\) in

\[
\frac{f(z)}{z-a}
\]

must lie inside the contour for the standard formula to contribute
\(2\pi i f(a)\).

%%%%%%%%%%%%%%%%%%%%%%%%%%%%%%%%%%%%%%%%%%%%%%%%%%%%%%%%%%%%
\subsubsection{Confusing Taylor and Laurent Series}

Taylor series contain only nonnegative powers.

Laurent series may contain negative powers and are used near isolated
singularities.

%%%%%%%%%%%%%%%%%%%%%%%%%%%%%%%%%%%%%%%%%%%%%%%%%%%%%%%%%%%%
\subsubsection{Confusing a Pole with an Essential Singularity}

A pole has only finitely many negative-power Laurent terms.

An essential singularity has infinitely many.

%%%%%%%%%%%%%%%%%%%%%%%%%%%%%%%%%%%%%%%%%%%%%%%%%%%%%%%%%%%%
\subsubsection{Computing the Wrong Laurent Coefficient}

The residue is specifically the coefficient of

\[
\boxed{
(z-a)^{-1}.
}
\]

%%%%%%%%%%%%%%%%%%%%%%%%%%%%%%%%%%%%%%%%%%%%%%%%%%%%%%%%%%%%
\subsubsection{Forgetting All Interior Singularities}

The Residue Theorem requires summing the residues of every isolated
singularity inside the contour.

%%%%%%%%%%%%%%%%%%%%%%%%%%%%%%%%%%%%%%%%%%%%%%%%%%%%%%%%%%%%
\subsection{Applications}
\label{subsec:complex-analysis-applications}
%%%%%%%%%%%%%%%%%%%%%%%%%%%%%%%%%%%%%%%%%%%%%%%%%%%%%%%%%%%%

\begin{applicationbox}
\textbf{Real Integrals.}
Contour integration and residues provide elegant methods for evaluating
many improper real integrals.

\medskip

\textbf{Differential Equations.}
Complex methods appear in transform techniques, Green's functions, and
the analysis of differential equations.

\medskip

\textbf{Fourier Analysis.}
Complex exponentials provide the natural language of Fourier series and
Fourier transforms.

\medskip

\textbf{Potential Theory.}
Real and imaginary parts of holomorphic functions are harmonic, connecting
complex analysis directly to Laplace's equation.

\medskip

\textbf{Fluid Mechanics.}
Two-dimensional incompressible and irrotational flows can be represented
using complex potentials.

\medskip

\textbf{Electrostatics.}
Conformal mappings can transform difficult two-dimensional boundary-value
problems into simpler geometries.

\medskip

\textbf{Signal Processing.}
Complex exponentials, transforms, poles, and residues are fundamental in
frequency-domain methods.

\medskip

\textbf{Control Theory.}
Poles and zeros of complex transfer functions determine stability and
system behavior.

\medskip

\textbf{Number Theory.}
Analytic continuation and complex-variable methods are central to analytic
number theory and functions such as the Riemann zeta function.

\medskip

\textbf{Geometry and Topology.}
Riemann surfaces, conformal geometry, winding numbers, and analytic
continuation connect complex analysis to geometry and topology.
\end{applicationbox}

%%%%%%%%%%%%%%%%%%%%%%%%%%%%%%%%%%%%%%%%%%%%%%%%%%%%%%%%%%%%
\subsection{Exercises}
\label{subsec:complex-analysis-exercises}
%%%%%%%%%%%%%%%%%%%%%%%%%%%%%%%%%%%%%%%%%%%%%%%%%%%%%%%%%%%%

\begin{exercise}[1]
Write

\[
z=1+i
\]

in polar and exponential form.
\end{exercise}

\begin{exercise}[1]
Compute

\[
|3-4i|.
\]
\end{exercise}

\begin{exercise}[1]
Find

\[
\overline{2+5i}.
\]
\end{exercise}

\begin{exercise}[1]
Find all cube roots of \(1\).
\end{exercise}

\begin{exercise}[1]
Use De Moivre's formula to compute

\[
(1+i)^8.
\]
\end{exercise}

\begin{exercise}[2]
Show that

\[
|zw|=|z||w|.
\]
\end{exercise}

\begin{exercise}[2]
Verify the triangle inequality for complex numbers using their
identification with vectors in \(\R^2\).
\end{exercise}

\begin{exercise}[2]
Determine whether

\[
f(z)=z^3
\]

satisfies the Cauchy--Riemann equations.
\end{exercise}

\begin{exercise}[2]
Determine where

\[
f(z)=|z|^2
\]

is complex differentiable.
\end{exercise}

\begin{exercise}[2]
Show that

\[
f(z)=\overline{z}
\]

is nowhere holomorphic.
\end{exercise}

\begin{exercise}[2]
Verify that the real and imaginary parts of

\[
f(z)=z^2
\]

are harmonic.
\end{exercise}

\begin{exercise}[2]
Find a harmonic conjugate of

\[
u(x,y)=x^2-y^2.
\]
\end{exercise}

\begin{exercise}[2]
Show that

\[
e^{z+2\pi i}=e^z.
\]
\end{exercise}

\begin{exercise}[2]
Find all values of

\[
\log(-1).
\]
\end{exercise}

\begin{exercise}[2]
Parameterize the circle

\[
|z|=R
\]

with positive orientation.
\end{exercise}

\begin{exercise}[2]
Evaluate

\[
\oint_{|z|=1}z^3\,dz.
\]
\end{exercise}

\begin{exercise}[2]
Evaluate

\[
\oint_{|z|=1}\frac1z\,dz.
\]
\end{exercise}

\begin{exercise}[3]
Use Cauchy's Integral Formula to evaluate

\[
\oint_{|z|=3}
\frac{\cos z}{z-1}\,dz.
\]
\end{exercise}

\begin{exercise}[3]
Evaluate

\[
\oint_{|z|=2}
\frac{e^z}{(z-1)^3}\,dz.
\]
\end{exercise}

\begin{exercise}[3]
Use Cauchy's derivative formula to express

\[
f^{(4)}(0)
\]

as a contour integral.
\end{exercise}

\begin{exercise}[3]
Use Liouville's Theorem to prove that a bounded polynomial must be
constant.
\end{exercise}

\begin{exercise}[3]
Use Liouville's Theorem to prove the Fundamental Theorem of Algebra.
\end{exercise}

\begin{exercise}[3]
Find the Laurent series of

\[
\frac{1}{z(1-z)}
\]

about \(z=0\) for

\[
0<|z|<1.
\]
\end{exercise}

\begin{exercise}[3]
Classify the singularity at \(z=0\) of

\[
\frac{\sin z}{z}.
\]
\end{exercise}

\begin{exercise}[3]
Classify the singularity at \(z=0\) of

\[
\frac1{z^4}.
\]
\end{exercise}

\begin{exercise}[3]
Classify the singularity at \(z=0\) of

\[
e^{1/z}.
\]
\end{exercise}

\begin{exercise}[3]
Compute

\[
\operatorname{Res}
\left(
\frac{1}{z(z-2)};
0
\right).
\]
\end{exercise}

\begin{exercise}[3]
Compute the residue of

\[
\frac{e^z}{(z-1)^2}
\]

at \(z=1\).
\end{exercise}

\begin{exercise}[3]
Use the Residue Theorem to evaluate

\[
\oint_{|z|=3}
\frac{1}{z(z-1)(z-2)}
\,dz.
\]
\end{exercise}

\begin{exercise}[3]
Use residues to evaluate

\[
\int_{-\infty}^{\infty}
\frac{dx}{x^2+a^2},
\qquad
a>0.
\]
\end{exercise}

\begin{exercise}[4]
Prove that the zeros of a nonzero holomorphic function are isolated.
\end{exercise}

\begin{exercise}[4]
Prove the Identity Theorem using the local power-series representation of
a holomorphic function.
\end{exercise}

\begin{exercise}[4]
Derive Cauchy's estimates from Cauchy's Integral Formula.
\end{exercise}

\begin{exercise}[4]
Prove the Maximum Modulus Principle.
\end{exercise}

\begin{exercise}[4]
Use Rouch\'e's Theorem to determine the number of zeros of

\[
z^5+10z+1
\]

inside a suitable circle.
\end{exercise}

\begin{exercise}[4]
Use the Argument Principle to count the zeros of a polynomial inside a
given contour.
\end{exercise}

\begin{exercise}[4]
Show that a M\"obius transformation maps generalized circles to
generalized circles.
\end{exercise}

\begin{exercise}[4]
Determine a M\"obius transformation mapping

\[
0,\quad1,\quad\infty
\]

to three prescribed distinct points.
\end{exercise}

\begin{exercise}[4]
Study the stereographic projection between the Riemann sphere and the
extended complex plane.
\end{exercise}

%%%%%%%%%%%%%%%%%%%%%%%%%%%%%%%%%%%%%%%%%%%%%%%%%%%%%%%%%%%%
\subsection{Conceptual Summary}
\label{subsec:complex-analysis-summary}
%%%%%%%%%%%%%%%%%%%%%%%%%%%%%%%%%%%%%%%%%%%%%%%%%%%%%%%%%%%%

Complex analysis studies functions

\[
\boxed{
f:\mathbb{C}\to\mathbb{C}.
}
\]

A complex number has the forms

\[
\boxed{
z=x+iy
}
\]

and

\[
\boxed{
z=re^{i\theta}.
}
\]

Euler's formula gives

\[
\boxed{
e^{i\theta}
=
\cos\theta+i\sin\theta.
}
\]

Complex differentiability is defined by

\[
\boxed{
f'(z)
=
\lim_{h\to0}
\frac{f(z+h)-f(z)}{h}.
}
\]

For

\[
f=u+iv,
\]

the Cauchy--Riemann equations are

\[
\boxed{
u_x=v_y,
\qquad
u_y=-v_x.
}
\]

Holomorphic functions are extraordinarily rigid:

\[
\boxed{
\text{holomorphic}
\Longrightarrow
\text{infinitely differentiable}
\Longrightarrow
\text{analytic}.
}
\]

Cauchy's Integral Theorem gives

\[
\boxed{
\oint_\gamma f(z)\,dz=0
}
\]

under appropriate holomorphicity and domain hypotheses.

Cauchy's Integral Formula gives

\[
\boxed{
f(a)
=
\frac{1}{2\pi i}
\oint_\gamma
\frac{f(z)}{z-a}\,dz.
}
\]

Laurent series extend Taylor series by allowing negative powers:

\[
\boxed{
f(z)
=
\sum_{n=-\infty}^{\infty}
c_n(z-a)^n.
}
\]

The residue is

\[
\boxed{
\operatorname{Res}(f;a)=c_{-1}.
}
\]

The Residue Theorem states

\[
\boxed{
\oint_\gamma f(z)\,dz
=
2\pi i
\sum_k
\operatorname{Res}(f;a_k).
}
\]

The Maximum Modulus Principle, Identity Theorem, Open Mapping Theorem, and
Liouville's Theorem demonstrate the rigidity of holomorphic functions.

Complex analysis connects algebra, geometry, topology, differential
equations, harmonic functions, Fourier analysis, physics, engineering,
and number theory through a remarkably unified theory.

%%%%%%%%%%%%%%%%%%%%%%%%%%%%%%%%%%%%%%%%%%%%%%%%%%%%%%%%%%%%
\subsection{Section Connections}
\label{subsec:complex-analysis-connections}
%%%%%%%%%%%%%%%%%%%%%%%%%%%%%%%%%%%%%%%%%%%%%%%%%%%%%%%%%%%%

\chapterconnection{
Complex Analysis extends calculus and real analysis from the real line to
the complex plane. Multivariable differentiation appears through the
Cauchy--Riemann equations, while topology enters through domains,
homotopy of contours, winding numbers, and branch behavior. Holomorphic
functions connect directly to harmonic functions and therefore to
Potential Theory. Cauchy's Integral Formula leads to power-series
representations, while Laurent series and residues provide tools for
studying singularities and evaluating integrals. Conformal mappings
connect the subject to geometry, fluid mechanics, and boundary-value
problems. Complex exponentials form the natural language of Fourier
Analysis, which is developed later in this chapter. Complex function
spaces and integral operators also provide important examples for
Functional Analysis and Operator Theory.
}

%%%%%%%%%%%%%%%%%%%%%%%%%%%%%%%%%%%%%%%%%%%%%%%%%%%%%%%%%%%%
% SECTION SOURCE TRACKING
%%%%%%%%%%%%%%%%%%%%%%%%%%%%%%%%%%%%%%%%%%%%%%%%%%%%%%%%%%%%

%------------------------------------------------------------
% 5.8 Complex Analysis
%------------------------------------------------------------
%
% Source basis:
%   - Standard undergraduate complex analysis.
%   - Standard introductory theory of holomorphic and meromorphic
%     functions.
%
% Used for:
%   - complex numbers and the complex plane;
%   - conjugation and modulus;
%   - argument and polar form;
%   - Euler's formula;
%   - De Moivre's formula and roots;
%   - complex functions, limits, and continuity;
%   - complex differentiation;
%   - holomorphic and entire functions;
%   - Cauchy--Riemann equations;
%   - harmonic functions and harmonic conjugates;
%   - complex exponential and trigonometric functions;
%   - logarithms, branches, and branch cuts;
%   - contour integration;
%   - ML estimate;
%   - antiderivatives and path independence;
%   - Cauchy's Integral Theorem;
%   - Cauchy's Integral Formula;
%   - Cauchy's derivative formula;
%   - analyticity and power series;
%   - Cauchy estimates;
%   - Liouville's Theorem;
%   - Fundamental Theorem of Algebra;
%   - Maximum Modulus Principle;
%   - Identity Theorem;
%   - zeros of holomorphic functions;
%   - Laurent series;
%   - isolated singularities;
%   - residues and Residue Theorem;
%   - real integral evaluation by residues;
%   - meromorphic functions;
%   - Argument Principle;
%   - winding numbers;
%   - Rouche's Theorem;
%   - Open Mapping Theorem;
%   - conformal mappings;
%   - Mobius transformations;
%   - Riemann sphere;
%   - Schwarz Lemma;
%   - analytic continuation.
%
% Notes:
%   - Real limits, continuity, differentiation, and integration are
%     developed earlier in Chapter 5.
%   - Multivariable calculus supplies the real-variable background
%     for the Cauchy--Riemann equations.
%   - Point-set and algebraic topology provide deeper foundations
%     for simply connected domains and winding numbers.
%   - Harmonic functions connect this section to Potential Theory.
%   - Complex exponentials and contour methods connect directly to
%     Fourier and Harmonic Analysis.
%
%%%%%%%%%%%%%%%%%%%%%%%%%%%%%%%%%%%%%%%%%%%%%%%%%%%%%%%%%%%%